\section{Prospective protocol and independent reconstruction}
\label{sec:protocol}

\subsection{Separate registrations, common logic}

The two studies share a target-domain and decision vocabulary but not a
numerical registration.  The \emph{same role} means the same continuum target,
input sector, map semantics, witness families, and causal, symmetry, and
resource categories.  Candidate law, lattice ladder, intervention cases,
numerical tolerances, and analytic convergence order are
realization-specific and were frozen separately.  A content seal binds
candidate and target code, state and scenario formulas,
maps, partitions, reducers, result schema, tests, dependencies, precision, and
stopping rule.  Registered execution remains disabled until the seal is
committed alone and identified by an annotated freeze tag.  A stage-opening
manifest is then committed alone.  The stage runs as one indivisible batch and
writes one no-clobber ledger, which is committed before inspection or replay.
Confirmation can open only if the committed development ledger passes both
its primary validator and the independent reconstruction.  This sequence is
shown at the foot of Fig.~\ref{fig:construction-map}.

The product development partition contains free light-mass, free heavy-mass,
upward mass-quench, and even-potential schedules.  Its confirmation partition
contains a downward mass quench and an odd-potential schedule.  The direct
development partition contains positive mass, negative mass, a sign-changing
quench, and a harmonic-$2$ potential.  Its confirmation partition contains a
three-stage mass pulse and a harmonic-$3$ potential.  Each partition uses four
formula-defined state families.  Product confirmation also includes a
predeclared cross-witness of its confirmation states against one development
schedule after the holdout is opened.

The confirmation schedules are single-use intervention holdouts, not draws
from a population model.  Because Eq.~\eqref{eq:certified-distance} is uniform
on $\Kin$, a finite witness-state split does not supply its logical force.  The
holdout tests whether the frozen map from a new exact intervention schedule to
the unchanged candidate continues to satisfy all uniform, observable,
structural, and replay gates.

\subsection{Frozen convergence gates}

For the product walk, every scenario must satisfy
\begin{equation}
 \begin{aligned}
 d_{256}&\le0.20,\qquad d_{512}\le0.10,\qquad d_{1024}\le0.05,\\
 d_{512}&\le0.60d_{256}+10^{-8},\qquad
 d_{1024}\le0.60d_{512}+10^{-8}.
 \end{aligned}
 \label{eq:product-absolute-gates}
\end{equation}
The scale was chosen before outcomes from the
first-order product estimate $T\Delta |km|$, with room for the registered
potential and certified tail.

The direct QCA uses its analytic finite-step signature.  Direct-only scenarios
must satisfy
\begin{equation}
 \begin{aligned}
 d_{320}&\le0.0040,\quad d_{640}\le0.0012,\quad d_{1280}\le0.0004,\\
 d_{640}&\le0.35d_{320}+10^{-10},\quad
 d_{1280}\le0.35d_{640}+10^{-10}.
 \end{aligned}
 \label{eq:direct-gates}
\end{equation}
whereas potential-split scenarios must satisfy
\begin{equation}
 \begin{aligned}
 d_{320}&\le0.020,\quad d_{640}\le0.011,\quad d_{1280}\le0.006,\\
 d_{640}&\le0.62d_{320}+10^{-10},\quad
 d_{1280}\le0.62d_{640}+10^{-10}.
 \end{aligned}
 \label{eq:potential-gates}
\end{equation}
An outcome-free 256-bit certificate bounded the complete direct one-step
remainder, telescoped it through $T=2$, charged the potential commutator,
target-cutoff tail, and numerical budget, and allowed potential-driven states
to leave the input band.  It used intermediate mode bands
$|r|\le13,17,20$ across the three refinements and charged the remaining Dyson
tail at every later step.  Its worst harmonic-$3$ envelopes were
$0.019662$, $0.007805$, and $0.002914$, below
Eq.~\eqref{eq:potential-gates}.  The certificate justified the gate scale; it did
not compute any registered trajectory, ratio, witness, or classification.

At the finest refinement, every product witness must have position and
momentum total-variation distance at most $0.05$, total- and local-current
error at most $0.10$, infidelity at most $0.0025$, and Zitterbewegung
current-series error at most $0.10$.  The corresponding direct gates are
$0.010$, $0.010$, $0.020$, $0.020$, $10^{-4}$, and $0.020$.
Both registrations require norm residual at most $10^{-12}$ and independent
state disagreement at most $10^{-11}$.  The direct potential scenarios also
specified candidate $\pi$-sector probability at most $10^{-24}$ in addition
to their analytic leakage tails; the empirical implementation of that
requirement is the exception documented in
Sec.~\ref{sec:post-result-audit}.

\subsection{Numerical and algorithmic separation}

Primary registered dynamics use \texttt{python-flint 0.9.0} Arb/Acb arithmetic
at 256 bits with outward enclosures; binary64 values are forbidden from the
registered dynamics.  The product primary evolves the candidate in real space
and the target in a retained Fourier basis.  Its independent route reconstructs
candidate gates and target evolution at 110 decimal digits with
\texttt{mpmath 1.3.0}.

For the direct update, separation is more explicit.  The primary
candidate applies the leading phase and recurrence in real space; the primary
target uses certified free blocks and spectral exponentials.  The independent
free candidate builds separate $2\times2$ momentum blocks.  Its potential
candidate constructs sparse phase, edge-scatter, and routing operators.  The
independent target rebuilds the harmonic generator and applies its exponential
through a separate high-precision route.  Candidate and verifier do not share
a candidate constructor, and target and verifier do not share a target
constructor.  This is an algorithmically separate reconstruction by the same
project, not an external replication.

\subsection{Infrastructure failures preserved rather than hidden}

The direct experiment required four separately sealed lifecycle lineages.  Version
0.1 stopped before ledger construction when the independent norm routine used
a real-symmetric eigensolver on a complex Hermitian Gram matrix.  No partial
ledger was written or inspected.  Version 0.2 wrote a primary-passing
development ledger, but replay found that a Fourier basis had been constructed
at ambient 53-bit precision before entry into the registered 256-bit context.
The certified upper bound was rigorous but failed the frozen $10^{-65}$ replay
tightness allowance.  Confirmation remained unopened.

Version 0.3 repaired precision ownership and independently reconstructed all
potential-free development scenarios, but its maximum-Gram-row-sum norm bound
became materially loose after the potential coupled registered momentum
columns.  It again stopped before confirmation.  Version 0.4 replaced only
this norm-bound reducer with a rigorous finite-precision similarity transform and
Gershgorin eigenvalue enclosure.  Approximate eigenvectors choose a basis but
carry no proof; the certified inverse, transform, and discs do.  Candidate,
target, maps, states, scenarios, partitions, thresholds, controls, and stopping
rules did not change across these infrastructure repairs.  The v0.2 and v0.3
primary classifications are not counted as verified results.  Their failure
records demonstrate why a second implementation must test certified
enclosures, not merely reproduce floating-point midpoints.

\subsection{Post-result semantic audit of the zone-edge reducer}
\label{sec:post-result-audit}

A hostile-reader audit after the v0.4 result found a shared factor-of-two error
in one reducer.  The specification defines $|q|\ge3\pi/4$ with
$q=2\pi r/N$, hence
\begin{equation}
 8|r|\ge3N.
 \label{eq:pi-sector-correct}
\end{equation}
The primary reducer, independent verifier, and non-gating rough-potential
control instead used $4|r|\ge3N$.  Since every discrete mode has
$|r|\le N/2$, that predicate selects an empty set.  All serialized empirical
$\pi$-probabilities and the rough-control value zero are therefore vacuous and
withdrawn.  The replay duplicated the predicate and is not independent evidence
for this metric.

The pre-outcome analytic leakage routine is separate and uses the correct
coarsest boundary $3N/8=120$.  Its serialized ledger values are approximately
$1.62\times10^{-261}$ for the harmonic-$2$ schedule and
$1.83\times10^{-164}$ for harmonic $3$, below the frozen $10^{-24}$ gate.
Thus the low-band leakage threshold remains supported analytically, while the
claim that its additional empirical reducer was executed correctly is
withdrawn.  No registered ledger is altered.  An audit-only recomputation of
the disclosed non-gating rough control with Eq.~\eqref{eq:pi-sector-correct}
gives $3.4860\times10^{-8}$, confirming the defect but neither repairing nor
broadening confirmation.  The serialized ledger classification field remains a
content-addressed historical fact; this manuscript reports it only with this
qualification.
